% ~1,5 páginas
\section{Abordagem proposta}
\label{sec:abordagem}

A \cref{fig:arquitetura} resume a arquitetura.

\begin{figure}[t]
\centering
\begin{tikzpicture}[
    node distance=4.5mm,
    caixa/.style={draw=corEscura, rounded corners=2pt, align=center,
      font=\scriptsize, inner sep=3pt, minimum height=9mm, text width=17mm},
    seta/.style={-{Stealth[length=1.6mm]}, draw=corEscura}]
  \node[caixa] (dados) {Cenas do WOMD\\(\un{20}{s}, \un{10}{Hz})};
  \node[caixa, right=of dados] (amostras) {Amostras\\(veículo, instante)};
  \node[caixa, right=of amostras] (prev)
    {Previsão\\vel.\ constante\\+ correção};
  \node[caixa, right=of prev] (plan) {Planeador A*\\$(s, v, t)$};
  \node[caixa, right=of plan] (aval) {Avaliação\\contra o registo};
  \draw[seta] (dados) -- (amostras);
  \draw[seta] (amostras) -- (prev);
  \draw[seta] (prev) -- (plan);
  \draw[seta] (plan) -- (aval);
\end{tikzpicture}
\caption{Arquitetura «prever para planear»: das cenas gravadas ao plano de
  velocidade do veículo autónomo.}
\label{fig:arquitetura}
\end{figure}

\subsection{Dados e amostras}

Usam-se dois ficheiros da versão de \un{20}{s} do WOMD, com 149 cenas.
Uma amostra é um par (veículo, instante). Em cada cena toma-se um
instante por segundo e, nele, todos os veículos a menos de \un{55}{m} do
veículo autónomo, em movimento (pelo menos \un{0,5}{m/s}) e observados sem
falhas no segundo anterior. O resultado são 13\,856 amostras de 2\,343
veículos. O alvo é o deslocamento do veículo 1, 3 e \un{5}{s} depois; o
de \un{5}{s} só existe em 10\,295 amostras, porque muitos veículos deixam
de ser observados antes disso, e as restantes não entram no treino nem na
avaliação desse horizonte.

\subsection{Características e modelos}

Cada amostra é descrita por 21 variáveis calculadas apenas com o passado e
o presente, em três grupos:
\begin{itemize}
  \item movimento do próprio veículo, no seu referencial: velocidade
    longitudinal e lateral, aceleração estimada, taxa de viragem e
    deslocamento no último segundo;
  \item veículo da frente: distância, velocidade e velocidade de
    aproximação, quando existe um veículo ou ciclista na mesma direção, a
    menos de \un{60}{m} e de \un{1,8}{m} de desvio lateral (acontece em
    66\% das amostras);
  \item relação com o veículo autónomo: posição e velocidade relativas,
    distância, velocidade de aproximação e tempo até à colisão.
\end{itemize}

Seguindo as duas ideias da \cref{sec:previsao}, o modelo não prevê a
posição futura em coordenadas do mapa. Prevê, no referencial do veículo,
a diferença entre o deslocamento real e o que a velocidade constante
daria:
\begin{equation}
  \hat{\mathbf{d}}_h = \mathbf{v}\,h + f_h(\mathbf{x}),
  \qquad h \in \{1, 3, 5\}~\mathrm{s},
\end{equation}
em que $\mathbf{v}$ é a velocidade atual, $\mathbf{x}$ o vetor de
características e $f_h$ um modelo treinado para o horizonte $h$. Assim, um
veículo que segue em frente tem correção nula seja qual for a direção da
estrada no mapa, e o modelo concentra-se nos desvios: travagens,
acelerações e curvas. Comparam-se três modelos para $f_h$: árvore de
decisão (profundidade 6), floresta aleatória (60 árvores, profundidade 8)
e \emph{gradient boosting} por histogramas (80 iterações).

\subsection{Planeador}

O planeador segue a decomposição caminho-velocidade
(\cref{sec:planeamento}): o caminho é dado e planeia-se a velocidade. Um
estado é $(s, v, t)$, a distância percorrida, a velocidade e o instante,
com passo de \un{0,1}{s} e horizonte de \un{3}{s}. As ações são quatro
acelerações: $-3$, $-1,5$, $0$ e \un{$+1,5$}{m/s$^2$}. O custo de um plano
soma o desconforto (quadrado da aceleração) com uma penalização elevada
por cada passo a menos de \un{3,5}{m} da posição prevista de um veículo; a
heurística do A* premeia o progresso. As posições previstas vêm de uma de
quatro fontes: nenhuma (os outros ficam onde estão), velocidade constante,
floresta aleatória ou o futuro real, que serve de limite superior.

Duas ressalvas sobre a implementação atual. A colisão entra como custo e
não como proibição, e a heurística não é admissível, pelo que o plano
devolvido não tem garantia de ser o ótimo. E o código, em Python, não
está otimizado: prever e planear demoram, cada um, cerca de \un{0,1}{s}
por instante, o que ainda não chega para tempo real.
